% =====================================================================
%  Slide preamble: classic Beamer theme (Warsaw), fonts, listings, TikZ
% =====================================================================

% ---------- Fonts and language ----------
% pdfLaTeX: standard Latin Modern. XeLaTeX/LuaLaTeX: Cambria Math if installed.
\usepackage{iftex}
\ifPDFTeX
  \usepackage[utf8]{inputenc}
  \usepackage[T1]{fontenc}
  \usepackage{lmodern}
  \usepackage{amsmath,amssymb}
\else
  \usefonttheme{professionalfonts}
  \usepackage{fontspec}
  \usepackage{amsmath,amssymb}
  \usepackage{unicode-math}
  \IfFontExistsTF{Cambria Math}
    {\setmathfont{Cambria Math}}
    {\setmathfont{Latin Modern Math}}
\fi
\usepackage[english]{babel}
\usepackage{csquotes}

% ---------- Classic theme ----------
\usetheme{Warsaw}
\setbeamertemplate{navigation symbols}{}
% Split footline as in the theme, with the frame number on the right
\setbeamertemplate{footline}{%
  \leavevmode%
  \hbox{%
  \begin{beamercolorbox}[wd=.5\paperwidth,ht=2.5ex,dp=1.125ex,leftskip=.3cm plus1fill,rightskip=.3cm]{author in head/foot}%
    \usebeamerfont{author in head/foot}\insertshortauthor
  \end{beamercolorbox}%
  \begin{beamercolorbox}[wd=.5\paperwidth,ht=2.5ex,dp=1.125ex,leftskip=.3cm,rightskip=.3cm]{title in head/foot}%
    \usebeamerfont{title in head/foot}\insertshorttitle\hfill\insertframenumber{} / \inserttotalframenumber
  \end{beamercolorbox}}%
  \vskip0pt}

% Outline at the start of each section
\AtBeginSection[]{%
  \begin{frame}<beamer>
    \tableofcontents[currentsection,hideothersubsections]
  \end{frame}}

% ---------- Packages ----------
\usepackage{booktabs,tabularx,array}
\usepackage{listings}
\usepackage{adjustbox}
\usepackage{tikz}
\usetikzlibrary{arrows.meta,positioning,fit,calc,shapes.geometric,backgrounds}
\usepackage[backend=biber,style=numeric-comp,sorting=none,maxbibnames=3,
            url=false,doi=false,eprint=false,isbn=false]{biblatex}
\addbibresource{bibliography/references.bib}
\renewcommand*{\bibfont}{\scriptsize}
\setbeamertemplate{bibliography item}{\insertbiblabel}
\newcolumntype{L}{>{\raggedright\arraybackslash}X}

% ---------- Listings (YANG, pyang trees, HTTP) ----------
\lstdefinelanguage{yang}{
  morekeywords={module,namespace,prefix,import,organization,contact,
    description,revision,reference,yang-version,identity,base,typedef,type,
    identityref,grouping,uses,container,list,key,leaf,leaf-list,config,
    augment,rpc,input,output,notification,leafref,path,default,units,
    must,when,enum,enumeration,mandatory},
  sensitive=true,
  morecomment=[l]{//},
  morecomment=[s]{/*}{*/},
  morestring=[b]"
}
\lstset{
  basicstyle=\ttfamily\scriptsize,
  keywordstyle=\bfseries,
  frame=single,
  columns=fullflexible, keepspaces=true,
  breaklines=true, upquote=true,
  aboveskip=4pt, belowskip=4pt
}

% ---------- Common TikZ styles (as in the paper) ----------
\tikzset{
  >=Stealth,
  box/.style={draw,rounded corners,minimum height=8mm,align=center,
              font=\small},
  lbl/.style={font=\scriptsize,midway,fill=white,inner sep=1pt},
  actor/.style={draw,rounded corners,fill=gray!8,minimum width=2.4cm,
                minimum height=7mm,align=center,font=\small},
  msg/.style={->,thick},
  ret/.style={->,dashed}
}

% Paper TikZ figure scaled to the slide: \slidefig[max height]{file}
\newcommand{\slidefig}[2][0.75\textheight]{%
  \begin{center}
    \begin{adjustbox}{max width=\textwidth,max totalheight=#1}
      \input{figures/#2}
    \end{adjustbox}
  \end{center}}

% Source note, bottom right
\newcommand{\source}[1]{\vfill\hfill{\tiny #1}}
